\documentclass[11pt]{article}

\usepackage[margin=1in]{geometry}
\usepackage{graphicx}
\usepackage{amsmath}
\usepackage{siunitx}
\usepackage{caption}
\usepackage{fancyhdr}
\usepackage{float}
\usepackage{hyperref}
\usepackage{enumitem}

\pagestyle{fancy}
\setlength{\headheight}{14pt}
\fancyhf{}
\lhead{Project \#4: Circuit Element Identification}
\rhead{\thepage}
\renewcommand{\headrulewidth}{0.4pt}

\title{\vspace{-2cm}Linear Circuits 2 \\ Circuit Element Identification Lab Report}
\author{Group Members: \\ Juvon Mondesir \\ Kenny Gallmon}
\date{}

\begin{document}

\maketitle
\thispagestyle{fancy}

\section*{Objectives}

We are given a concealed circuit with the following specifications:

\begin{itemize}[nosep]
    \item The concealed circuit contains two circuit elements only:
    Element X and Element Y
    \item The elements are connected in series
    \item Both elements X and Y can be a resistor, inductor, or capacitor
    \item The concealed circuit must contain at least one resistor
\end{itemize}

\begin{figure}[H]
    \centering
    \includegraphics[width=0.85\textwidth]{figures/fig01_concealed_circuit.png}
    \caption{Concealed circuit under test. Node A is connected to a red
    wire, Node B to a black wire, and Node C to a green wire.}
    \label{fig:concealed_circuit}
\end{figure}

We are given four tasks:

\begin{itemize}[nosep]
    \item Identify the type of element that is connected between node A
    and node C.
    \item Determine the value for element X. The predicted value must fall
    within $\pm 20\%$ of the actual component value.
    \item Identify the type of element that is connected between node B
    and node C.
    \item Determine the value for element Y. The predicted value must fall
    within $\pm 20\%$ of the actual component value.
\end{itemize}

\section*{Dissection of Experimental Plan: Part One}

To begin with, the first part of our plan was to determine the combination
of elements X and Y. We concluded that there were only three combinations,
which include RR, RL, and RC. The first test that we performed was to apply
a square wave input to observe the phase difference of the waveforms across
both elements X and Y.

\section*{Experimental Results: Part One}

\begin{figure}[H]
    \centering
    \includegraphics[width=0.55\textwidth]{figures/fig02_funcgen_settings.png}
    \caption{Function generator settings: $1$\,kHz, $10$\,V$_{pp}$ square
    wave}
\end{figure}

\begin{figure}[H]
    \centering
    \includegraphics[width=0.95\textwidth]{figures/fig03_scope_waveform_p1.png}
    \caption{Oscilloscope output of the waveforms across component X and Y}
\end{figure}

Using a 1\,kHz, 10\,V$_{pp}$ square wave, we observe the oscilloscope output
of component X and Y. Utilizing the oscilloscope, we find that there is no
phase difference. As a result, we determine that the concealed circuit has
two resistors in series.

\section*{Dissection of Experimental Plan: Part Two}

After gaining our Part One results, we determined the concealed circuit was
a voltage divider. We had $V_{ab}$ as the input voltage, $V_{cb}$ as the
voltage drop across resistor Y, and $V_{ac}$ as the voltage drop across
resistor X, equal to $V_{ab} - V_{cb}$.

\begin{figure}[H]
    \centering
    \includegraphics[width=0.85\textwidth]{figures/fig04_voltage_divider_diagram.png}
    \caption{Concealed circuit modeled as a voltage divider, with Element X
    and Element Y both resistive}
\end{figure}

We required two equations to determine the values for both resistors: an
equation for the voltage drop across resistors in series, and an equation
for the equivalent resistance between two series resistors.

\begin{figure}[H]
    \centering
    \includegraphics[width=0.32\textwidth]{figures/fig_eq01_vout_formula.png}
    \vspace{4pt}

    \includegraphics[width=0.28\textwidth]{figures/fig_eq02_rxry_formula.png}
    \caption{Voltage-divider equation and series-resistance equation}
\end{figure}

\section*{Experimental Results: Part Two}

The second test performed on the concealed circuit was a current and
voltage measurement.

\begin{figure}[H]
    \centering
    \includegraphics[width=0.55\textwidth]{figures/fig05_scope_current.png}
    \caption{Measured AC current, $I = 0.2204$\,mA}
\end{figure}

We measure an AC current $I$ of $0.22$\,mA.

\begin{figure}[H]
    \centering
    \includegraphics[width=0.95\textwidth]{figures/fig06_scope_waveform_p2.png}
    \caption{Oscilloscope voltage and current measurement across the
    concealed circuit}
\end{figure}

With $V_{ab} = 5$\,V and $V_{cb} = 2$\,V, we substitute the values into the
equation for equivalent resistance between two series resistors:

\begin{figure}[H]
    \centering
    \includegraphics[width=0.5\textwidth]{figures/fig_eq03_rxry_calc.png}

    \vspace{6pt}
    \includegraphics[width=0.22\textwidth]{figures/fig_eq04_ry_ratio.png}

    \vspace{6pt}
    \includegraphics[width=0.55\textwidth]{figures/fig_eq05_final_values.png}
    \caption{Solving for $R_x$ and $R_y$ from the measured current and
    voltages}
\end{figure}

We concluded that Element X was a 9\,k$\Omega$ resistor and that Element Y
was a 14\,k$\Omega$ resistor.

\begin{figure}[H]
    \centering
    \includegraphics[width=0.55\textwidth]{figures/fig07_handcalc.png}
    \caption{Hand calculations backing the previous results}
\end{figure}

\section*{Conclusion}

We observed that both the input and output signals had no phase difference,
by comparing the square wave behavior of $V_{ab}$ and $V_{cb}$. This alerted
us to the conclusion that the concealed circuit was an R--R combination.
Once we realized it was two resistors in series, we measured the input and
output voltages, took note of the current flowing through the concealed
circuit, and used the voltage divider and series resistance equations to
determine the value of each resistor. Overall, the results aligned with our
expectations, confirming the approach for identifying and solving the
concealed circuit worked.

\end{document}